% ---
% titre: Amplification et simplification de fractions
% theme: NO
% chapitre: Nombres rationnels
% sous_chapitre: Amplifier et simplifier
% niveau: [10e]
% type: EX
% tags: [c-fractions, c-livrets, c-multiplication, c-division, f-individuel, p-approfondissement]
% difficulte: 2
% source: latex
% ---

\header{Droite numérique : placer, comparer, encadrer}


\def\x{1.3cm}
\def\arraystretch{2.2}
\vfill
{\bfseries\color{themeColor}Exercice 1.} Amplifie la fraction pour obtenir ce qui est demandé.

\begin{enumerate}[leftmargin=*,align=parleft, label=\alph*), itemsep=9pt]
\item $\dfrac{3}{4}$ pour obtenir le dénominateur $60$ \quad $\Longrightarrow$
\item $\dfrac{3}{4}$ pour obtenir le numérateur $36$ \quad $\Longrightarrow$
\item $\dfrac{2}{5}$ pour obtenir le dénominateur $35$ \quad $\Longrightarrow$
\item $\dfrac{7}{6}$ pour obtenir le numérateur $84$ \quad $\Longrightarrow$
\end{enumerate}

\vfill
{\bfseries\color{themeColor}Exercice 2.} Écris toutes les fractions égales à $\dfrac{3}{7}$, obtenues en amplifiant, dont le dénominateur est compris entre $10$ et $50$.

\vspace{100pt}

\vfill
{\bfseries\color{themeColor}Exercice 3.} Trouve la fraction!

\begin{enumerate}[leftmargin=*,align=parleft, label=\alph*), itemsep=50pt]
\item On l'obtient en amplifiant $\dfrac{5}{6}$ et la somme de son numérateur et de son dénominateur vaut $44$.
\item On l'obtient en amplifiant $\dfrac{3}{8}$ et la somme de son numérateur et de son dénominateur vaut $55$.
\end{enumerate} 

%\bl[\linewidth]
\vfill
\newpage


{\bfseries\color{themeColor}Exercice 4.} Simplifie la fraction pour obtenir ce qui est demandé.
 
\begin{enumerate}[leftmargin=*,align=parleft, label=\alph*), itemsep=9pt]
\item $\dfrac{45}{60}$ pour obtenir le dénominateur $4$ \quad $\Longrightarrow$
\item $\dfrac{36}{48}$ pour obtenir le numérateur $3$ \quad $\Longrightarrow$
\item $\dfrac{14}{35}$ pour obtenir le dénominateur $5$ \quad $\Longrightarrow$
\item $\dfrac{84}{72}$ pour obtenir le numérateur $7$ \quad $\Longrightarrow$
\end{enumerate}

\vfill
{\bfseries\color{themeColor}Exercice 5.} Écris toutes les fractions égales à $\dfrac{36}{60}$, obtenues en simplifiant par un diviseur commun de $36$ et $60$ (autre que $1$).

\vspace{100pt}

\vfill
{\bfseries\color{themeColor}Exercice 6.} Léa écrit : $\dfrac{6}{9}=\dfrac{3}{6}$, car \og j'ai enlevé $3$ en haut et en bas\fg{}.
 
\begin{enumerate}[leftmargin=*,align=parleft, label=\alph*), itemsep=20pt]
\item Explique pourquoi c'est faux.\\
\bl[\linewidth]\\
\bl[\linewidth]\\
\bl[\linewidth]\\
\bl[\linewidth]
\item Écris une fraction égale à $\dfrac{6}{9}$ de dénominateur $3$.
\end{enumerate}

\vfill
\newpage